\subsection{一元多项式在域上的整除性}

交换域 K 上一个未定元的多项式环 K{[}X{]} 是主理想整环（IV，第 11 页，命题 11）。其分式域是系数取自 K 的、关于未定元 X 的有理函数域 K(X)。环 K{[}X{]} 包含次数为 0 的多项式构成的子环，即常数域，它与 K 等同；K* 中的元素在 K 中可逆，从而在 K{[}X{]} 中也可逆；反之，公式 \(\deg(uv) = \deg(u) + \deg(v)\) 表明每个可逆多项式都有次数 0；K{[}X{]} 的可逆元素群 U 因此恰好是 K*。于是两个相伴多项式仅相差一个非零常数因子；特别地，每个相伴多项式类包含唯一一个首一多项式。乘法群 K(X)* 的由首一多项式生成的子群因此从每个相伴有理函数类中恰好包含一个元素，从而与群

\[
\mathcal {P} ^ {*} = \mathrm{K} (\mathrm{X}) ^ {*} / \mathrm{U}
\]

\(\mathrm{K}(\mathrm{X})\) 的主分式理想组成的群。特别地，每当提及域 \(\mathrm{K}(\mathrm{X})\) 中（关于环 \(\mathrm{K}[\mathrm{X}]\) ）的 gcd 或 lcm 时，通常理解为这些是首一多项式（或 0）的商；这一约定使我们能够谈论一族有理函数的 gcd 或 lcm。

K{[}X{]} 的不可约元素正是通常意义下的不可约多项式（IV，第 13 页，定义 2），而首一不可约多项式的集合构成了 K{[}X{]} 的不可约元素的一个代表系.

一次多项式总是不可约的。若 K 是代数闭域，则其逆命题成立（V，第 19 页，命题 1）；因此在这种情况下，每个多项式 \(p(\mathbf{X})\) （ \(n\) 次的，属于 \(\mathbf{K}[\mathbf{X}]\) ）都可以唯一地（不计因子的顺序）写成如下形式

\[
p (\mathbf {X}) = c \left(\mathbf {X} - a _ {1}\right) \left(\mathbf {X} - a _ {2}\right) \dots \left(\mathbf {X} - a _ {n}\right)
\]

，其中 \(c\) 与诸 \(a_{i}\) 都是 \(\mathbf{K}\) 的元素.

命题6.——对于任意域 K，K{[}X{]} 中首一不可约多项式的集合是无限的。

事实上，给定任意非空有限族 \((p_{i})\) ( \(1 \leqslant i \leqslant n\) ) 由不同的首一不可约多项式组成，多项式 \(\left(\prod_{i=1}^{n} p_{i}\right) + 1\) 不可逆，且该多项式的任何首一不可约因子 q 必然与所有 \(p_{i}\) 不同，否则它会整除 1。

